\documentclass[aspectratio=169,11pt]{beamer}
\usetheme{Madrid}
\usecolortheme{dolphin}
\setbeamertemplate{navigation symbols}{}
\setbeamertemplate{caption}[numbered]

\usepackage[utf8]{inputenc}
\usepackage[T1]{fontenc}
\usepackage[spanish,es-noshorthands]{babel}
\usepackage{amsmath,amssymb,amsthm,mathtools}
\usepackage{tikz}
\usepackage{pgfplots}
\pgfplotsset{compat=1.18}
\usepackage{booktabs}
\usepackage{array}

\title[Prob. y Estadística]{Clase 17: Ley de los grandes números}
\subtitle{Unidad 4: Función generadora de momentos}
\institute[UNS]{Universidad Nacional del Sur \\ Departamento de Matemática}
\author{Probabilidad y Estadística (5765)}
\date{}

\begin{document}

\begin{frame}
\titlepage
\end{frame}

\begin{frame}{Contenidos}
\tableofcontents
\end{frame}

\section{Motivación}

\begin{frame}{La intuición frecuencial de la probabilidad}
\begin{block}{Pregunta fundamental}
Cuando decimos que ``la probabilidad de cara al tirar una moneda equilibrada es $0.5$'', ¿qué queremos decir exactamente? Intuitivamente: si tiramos la moneda muchas veces, la \textbf{proporción} de caras se acerca a $0.5$.
\end{block}

\begin{alertblock}{Objetivo de esta clase}
Formalizar y demostrar rigurosamente esta intuición: la \textbf{ley de los grandes números}, que justifica matemáticamente la interpretación frecuencial de la probabilidad, y es la base teórica de la simulación y la inferencia estadística.
\end{alertblock}
\end{frame}

\section{Convergencia en probabilidad}

\begin{frame}{Sucesiones de variables aleatorias}
\begin{block}{Contexto}
Consideramos una sucesión de variables aleatorias $X_1, X_2, X_3, \dots$ y queremos dar sentido a la afirmación ``$X_n$ se acerca a un valor (o variable) $X$ cuando $n\to\infty$''. A diferencia de sucesiones numéricas, esto requiere una definición probabilística, porque cada $X_n$ es aleatoria.
\end{block}

\begin{definition}[Convergencia en probabilidad]
Decimos que la sucesión $\{X_n\}$ \textbf{converge en probabilidad} a una constante (o variable aleatoria) $X$, y escribimos $X_n \xrightarrow{P} X$, si para todo $\varepsilon > 0$:
\[
\lim_{n\to\infty} P\big(|X_n - X| \geq \varepsilon\big) = 0.
\]
\end{definition}
\end{frame}

\begin{frame}{Interpretación}
\begin{block}{Lectura}
Para cualquier margen de error $\varepsilon$ que fijemos (por pequeño que sea), la probabilidad de que $X_n$ se aparte de $X$ más de $\varepsilon$ se vuelve arbitrariamente chica a medida que $n$ crece.
\end{block}

\begin{alertblock}{Distinción importante}
La convergencia en probabilidad \textbf{no} dice que $X_n(\omega) \to X(\omega)$ para cada resultado particular $\omega$ del espacio muestral (eso sería convergencia ``casi segura'', un concepto más fuerte que no se estudia en este curso). Es una afirmación sobre cómo se comporta la \textbf{distribución} de la probabilidad de discrepancia a medida que $n$ crece.
\end{alertblock}
\end{frame}

\section{Ley débil de los grandes números}

\begin{frame}{Planteo del problema}
\begin{block}{Escenario}
Sean $X_1, X_2, \dots, X_n$ variables aleatorias \textbf{independientes e idénticamente distribuidas} (i.i.d.), con $E(X_i) = \mu$ y $\text{Var}(X_i) = \sigma^2 < \infty$ para todo $i$ (por ejemplo: resultados de tiradas repetidas de un dado, o mediciones repetidas de un instrumento).
\end{block}

Definimos el \textbf{promedio muestral}:
\[
\overline{X}_n = \frac{X_1 + X_2 + \cdots + X_n}{n} = \frac{1}{n}\sum_{i=1}^n X_i.
\]

\begin{alertblock}{Pregunta}
¿Qué le pasa a $\overline{X}_n$ a medida que $n \to \infty$?
\end{alertblock}
\end{frame}

\begin{frame}{Esperanza y varianza del promedio muestral}
\begin{block}{Cálculo previo}
Usando linealidad de la esperanza (Clase 14):
\[
E(\overline{X}_n) = E\left(\frac{1}{n}\sum_{i=1}^n X_i\right) = \frac{1}{n}\sum_{i=1}^n E(X_i) = \frac{1}{n}(n\mu) = \mu.
\]
Usando que los $X_i$ son independientes (Clase 15, propiedad de varianza de suma de independientes):
\[
\text{Var}(\overline{X}_n) = \text{Var}\left(\frac{1}{n}\sum_{i=1}^n X_i\right) = \frac{1}{n^2}\sum_{i=1}^n \text{Var}(X_i) = \frac{1}{n^2}(n\sigma^2) = \frac{\sigma^2}{n}.
\]
\end{block}

\begin{alertblock}{Observación crucial}
$E(\overline{X}_n) = \mu$ para todo $n$ (el promedio muestral es siempre ``centrado'' en $\mu$), pero $\text{Var}(\overline{X}_n) = \sigma^2/n \to 0$ cuando $n\to\infty$: la dispersión del promedio se reduce a medida que promediamos más datos.
\end{alertblock}
\end{frame}

\begin{frame}{Enunciado de la ley débil}
\begin{theorem}[Ley débil de los grandes números]
Sean $X_1,\dots,X_n$ i.i.d.\ con $E(X_i)=\mu$ y $\text{Var}(X_i)=\sigma^2<\infty$. Entonces el promedio muestral converge en probabilidad a $\mu$:
\[
\overline{X}_n \xrightarrow{P} \mu, \qquad \text{es decir,} \qquad \lim_{n\to\infty} P\big(|\overline{X}_n - \mu| \geq \varepsilon\big) = 0 \quad \forall \varepsilon>0.
\]
\end{theorem}

\begin{block}{Nombre}
Se llama ``débil'' porque solo garantiza convergencia en probabilidad; existe una versión ``fuerte'' (con convergencia casi segura) que requiere herramientas más avanzadas y queda fuera del alcance de este curso.
\end{block}
\end{frame}

\begin{frame}{Demostración usando la desigualdad de Chebyshev}
\begin{proof}
Aplicamos la desigualdad de Chebyshev (Clase 15) a la variable $\overline{X}_n$, que tiene esperanza $\mu$ y varianza $\sigma^2/n$. Para cualquier $\varepsilon>0$:
\[
P\big(|\overline{X}_n - \mu| \geq \varepsilon\big) \leq \frac{\text{Var}(\overline{X}_n)}{\varepsilon^2} = \frac{\sigma^2/n}{\varepsilon^2} = \frac{\sigma^2}{n\,\varepsilon^2}.
\]
Como $\sigma^2$ y $\varepsilon$ están fijos, al tomar $n\to\infty$ el lado derecho tiende a $0$:
\[
\lim_{n\to\infty} P\big(|\overline{X}_n-\mu|\geq\varepsilon\big) \leq \lim_{n\to\infty} \frac{\sigma^2}{n\varepsilon^2} = 0.
\]
Como una probabilidad no puede ser negativa, concluimos $\lim_{n\to\infty}P(|\overline{X}_n-\mu|\geq\varepsilon)=0$, es decir, $\overline{X}_n \xrightarrow{P} \mu$. $\blacksquare$
\end{proof}
\end{frame}

\begin{frame}{Ejemplo numérico: acotando con Chebyshev}
\begin{example}
Se mide el tiempo de reacción de personas a un estímulo. Se sabe que cada medición tiene $\sigma = 0.5$ segundos de desviación estándar (distribución desconocida). ¿Cuántas mediciones independientes $n$ se necesitan para garantizar, usando Chebyshev, que
\[
P\big(|\overline{X}_n - \mu| \geq 0.1\big) \leq 0.05\, ?
\]
\end{example}

\medskip
Por Chebyshev, $P(|\overline{X}_n-\mu|\geq 0.1) \leq \dfrac{\sigma^2}{n(0.1)^2} = \dfrac{0.25}{0.01\, n} = \dfrac{25}{n}$.

Queremos $\dfrac{25}{n} \leq 0.05 \iff n \geq \dfrac{25}{0.05} = 500$.

\textbf{Se necesitan al menos $500$ mediciones} para esa garantía (cota conservadora; en la práctica, con el TCL de la Clase 18 se obtienen cotas más ajustadas).
\end{frame}

\section{Interpretación frecuencial de la probabilidad}

\begin{frame}{Caso particular: frecuencia relativa}
\begin{block}{Aplicación clave}
Sea $A$ un evento con $P(A) = p$, y repitamos el experimento $n$ veces de forma independiente. Definimos $X_i = 1$ si $A$ ocurre en el ensayo $i$, y $X_i=0$ en caso contrario, es decir $X_i \sim \text{Bernoulli}(p)$ i.i.d.
\end{block}

Entonces $\overline{X}_n = \dfrac{1}{n}\sum_{i=1}^n X_i$ es exactamente la \textbf{frecuencia relativa} de ocurrencia de $A$ en las $n$ repeticiones (proporción de éxitos).

\begin{alertblock}{Consecuencia}
Como $E(X_i)=p$ y $\text{Var}(X_i)=p(1-p)$, la ley débil da
\[
\overline{X}_n \xrightarrow{P} p,
\]
es decir: \textbf{la frecuencia relativa de un evento converge en probabilidad a su probabilidad teórica}. Esto justifica formalmente la interpretación frecuencial (empírica) de la probabilidad usada informalmente desde el comienzo del curso.
\end{alertblock}
\end{frame}

\begin{frame}{Ejemplo: simulación conceptual de una moneda}
\begin{example}
Se tira una moneda equilibrada ($p=0.5$) $n$ veces. Sea $\hat{p}_n$ la frecuencia relativa de caras. Por la ley débil, $\hat p_n \xrightarrow{P} 0.5$.

\medskip
Usando Chebyshev con $\sigma^2 = p(1-p) = 0.25$: para $\varepsilon = 0.05$,
\[
P(|\hat p_n - 0.5| \geq 0.05) \leq \frac{0.25}{n(0.05)^2} = \frac{0.25}{0.0025\, n} = \frac{100}{n}.
\]
\begin{itemize}
\item Con $n=1000$: la cota es $100/1000 = 0.10$.
\item Con $n=10000$: la cota es $100/10000 = 0.01$.
\end{itemize}
Se observa cómo, al aumentar el número de tiradas, la probabilidad de que la frecuencia relativa se aleje de $0.5$ más de $5$ puntos porcentuales se hace cada vez más chica —consistente con lo que uno observaría si efectivamente simulara estas tiradas.
\end{example}
\end{frame}

\begin{frame}{Ejemplo: control de calidad}
\begin{example}
Una fábrica produce artículos con una proporción $p=0.02$ de defectuosos (desconocida en la práctica, pero supongamos que es el valor real). Se inspecciona una muestra de $n=2000$ artículos independientes. Sea $\hat p$ la proporción muestral de defectuosos.

\medskip
$E(\hat p) = 0.02$, $\text{Var}(\hat p) = \dfrac{p(1-p)}{n} = \dfrac{0.02(0.98)}{2000} = 0.0000098$.

Por Chebyshev, con $\varepsilon = 0.01$:
\[
P(|\hat p - 0.02| \geq 0.01) \leq \frac{0.0000098}{(0.01)^2} = \frac{0.0000098}{0.0001} = 0.098.
\]
Es decir, hay a lo sumo un $9.8\%$ de probabilidad de que la proporción muestral se aparte más de un punto porcentual de la proporción real de defectuosos. (En la Clase 18 refinaremos esta estimación usando el TCL.)
\end{example}
\end{frame}

\section{Visualización e implicancias}

\begin{frame}{Visualización: la varianza del promedio se contrae}
\begin{block}{Idea gráfica}
A medida que $n$ crece, la distribución de $\overline{X}_n$ se concentra cada vez más alrededor de $\mu$, porque $\text{Var}(\overline{X}_n) = \sigma^2/n \to 0$. El siguiente gráfico esquematiza cómo se ``achica'' la dispersión de $\overline X_n$ (en el eje vertical, una densidad esquemática) a medida que aumenta $n$.
\end{block}

\begin{center}
\begin{tikzpicture}[scale=0.85]
\begin{axis}[
  width=10cm, height=5cm,
  xlabel={valores de $\overline{X}_n$}, ylabel={densidad (esquemática)},
  xmin=-4, xmax=4, ymin=0, ymax=1.6,
  legend pos=north east, legend style={font=\tiny},
  ticks=none
]
\addplot[domain=-4:4, samples=100, thick, blue] {1/sqrt(2*pi*1.2)*exp(-x^2/(2*1.2))};
\addplot[domain=-4:4, samples=100, thick, red] {1/sqrt(2*pi*0.4)*exp(-x^2/(2*0.4))};
\addplot[domain=-4:4, samples=100, thick, green!60!black] {1/sqrt(2*pi*0.1)*exp(-x^2/(2*0.1))};
\legend{$n$ chico, $n$ mediano, $n$ grande}
\end{axis}
\end{tikzpicture}
\end{center}
\end{frame}

\begin{frame}{Ejemplo: promedio de calificaciones}
\begin{example}
Las calificaciones de un examen (escala $0$ a $100$) tienen $\mu=68$ y $\sigma=15$ (distribución desconocida). Se promedian los exámenes de $n=225$ estudiantes elegidos al azar e independientes entre sí. Acotar, con Chebyshev, la probabilidad de que el promedio muestral difiera de $68$ en más de $3$ puntos.
\end{example}

\medskip
$\text{Var}(\overline X_{225}) = \dfrac{15^2}{225} = \dfrac{225}{225} = 1$. Por Chebyshev con $k=3$:
\[
P(|\overline X_{225}-68|\geq 3) \leq \frac{1}{3^2} = \frac{1}{9} \approx 0.111.
\]
Con $225$ exámenes, la probabilidad de que el promedio se aparte más de $3$ puntos de la media poblacional es a lo sumo $11.1\%$: el promedio muestral es un estimador cada vez más preciso de $\mu$ a medida que $n$ crece.
\end{frame}

\begin{frame}{Alcance y limitaciones de la ley débil}
\begin{alertblock}{¿Qué garantiza y qué no?}
\begin{itemize}
\item Garantiza que, para $n$ grande, es \textbf{muy probable} que $\overline X_n$ esté cerca de $\mu$.
\item No garantiza nada sobre una realización \emph{particular}: en una muestra concreta, $\overline X_n$ puede alejarse de $\mu$ (solo se vuelve cada vez menos probable).
\item Requiere que $\text{Var}(X_i)$ sea finita; existen versiones más generales (fuera del alcance de este curso) que solo requieren $E(|X_i|)<\infty$.
\item Es la base teórica de: simulación de Monte Carlo, estimación de probabilidades mediante frecuencias relativas, y de gran parte de la inferencia estadística que se estudiará en unidades posteriores.
\end{itemize}
\end{alertblock}
\end{frame}

\begin{frame}{Ejemplo: encuesta de intención de voto}
\begin{example}
En un sondeo, se consulta a $n=1500$ votantes independientes si apoyan a un candidato. Supongamos que la proporción real de apoyo es $p=0.48$ (desconocida para quien realiza la encuesta). Acotar, con Chebyshev, la probabilidad de que la proporción muestral $\hat p$ se aparte de $0.48$ en más de $0.03$ (tres puntos porcentuales).
\end{example}

\medskip
$\text{Var}(\hat p) = \dfrac{p(1-p)}{n} = \dfrac{0.48(0.52)}{1500} = \dfrac{0.2496}{1500} \approx 0.0001664$.
\[
P(|\hat p - 0.48|\geq 0.03) \leq \frac{0.0001664}{0.03^2} = \frac{0.0001664}{0.0009} \approx 0.1849.
\]
La cota de Chebyshev da a lo sumo $18.5\%$ de probabilidad de error mayor a $3$ puntos porcentuales. Este tipo de cálculo es la base del llamado ``margen de error'' que se reporta en las encuestas (con el TCL, Clase 18, se obtienen cotas bastante más ajustadas que esta).
\end{frame}

\section{Resumen}

\begin{frame}{Resumen}
\begin{block}{Conceptos clave}
\begin{itemize}
\item \textbf{Convergencia en probabilidad}: $X_n \xrightarrow{P} X$ si $P(|X_n-X|\geq\varepsilon)\to 0$ para todo $\varepsilon>0$.
\item \textbf{Ley débil de los grandes números}: si $X_1,\dots,X_n$ son i.i.d.\ con media $\mu$ y varianza finita $\sigma^2$, entonces $\overline{X}_n \xrightarrow{P} \mu$.
\item \textbf{Demostración}: se obtiene directamente de la desigualdad de Chebyshev aplicada a $\overline{X}_n$, usando que $E(\overline{X}_n)=\mu$ y $\text{Var}(\overline{X}_n)=\sigma^2/n \to 0$.
\item \textbf{Caso particular}: la frecuencia relativa de un evento converge en probabilidad a su probabilidad teórica —justificación formal de la interpretación frecuencial.
\end{itemize}
\end{block}

\begin{alertblock}{Próxima clase}
El \textbf{teorema central del límite} refina este resultado: no solo dice que $\overline{X}_n$ se acerca a $\mu$, sino que describe \emph{cómo} se distribuyen las fluctuaciones alrededor de $\mu$ (aproximadamente normales), lo que permite calcular probabilidades aproximadas de manera mucho más precisa que con Chebyshev.
\end{alertblock}
\end{frame}

\end{document}
